% !TeX root = ../main.tex
% =========================================================================
% 9. PROYECTOS DE INVESTIGACIÓN VINCULADOS
% Basado formalmente en Doc01b (Solicitud de Adhesión institucional)
% =========================================================================
\section{Proyectos de investigación vinculados}

El presente trabajo de titulación, desarrollado por los estudiantes \textbf{\autorUnoProyecto} y \textbf{\autorDosProyecto}, se ejecuta bajo la dirección y tutoría del docente investigador Ing. Joe Llerena Izquierdo, M.Sc., Ph.D., y se encuentra formalmente articulado a las actividades científicas del \textbf{Grupo de Investigación en Enseñanza-Aprendizaje de las Ciencias para la Ingeniería (GIEACI)} de la Universidad Politécnica Salesiana, Sede Guayaquil.

Específicamente, los postulantes solicitan su adhesión formal al proyecto de investigación institucional titulado:
\begin{center}
    \textbf{\enquote{Percepciones y dilemas éticos de estudiantes de Ingeniería frente al uso de IA generativa en el aprendizaje}}
\end{center}
el cual se inscribe en la línea de investigación institucional de la carrera de Computación denominada: \textit{\enquote{Integración de Tecnologías Emergentes en la Educación en Ingeniería y Sistemas Inteligentes}}.

La vinculación con dicho proyecto de investigación se fundamenta en que el desarrollo del agente conversacional no solo aborda la dimensión técnica del software y la ingeniería de datos sobre SQLite, sino que aporta evidencia empírica directa sobre los dilemas de confiabilidad, sesgo y ética algorítmica derivados del uso de modelos generativos fundacionales en entornos de interacción humana real \parencite{educsci14121330,Ayala-Carabajo2026}. Los hallazgos sobre alucinaciones de precios, gobernanza de datos personales bajo la LOPDP ecuatoriana y patrones de interacción en mensajería instantánea enriquecerán las publicaciones y análisis del grupo GIEACI respecto a la adopción responsable de la inteligencia artificial.
